\section{From an economic event to permission to post}
\label{sec:posting-permission}

A delivery can change revenue, inventory and receivables at once. Before any
book changes, the method needs to know which contract was performed, which
opening state it consumes, which evidence supports that reading and which
accounting profile applies. Balance alone cannot supply those facts. AccDSL
therefore separates the economic basis, permission to use a treatment, proposed
entries and the transaction that makes them durable.

The event basis binds an immutable event to opening state, context, branch,
time view, chart, authority, method, evidence, admissions, assumptions and view.
Its business identity is the four-tuple of entity, book, branch and business
event. The branch declares whether it is synthetic or real. A synthetic teaching
profile can justify a teaching posting; it cannot authorize a real book merely
because the resulting arithmetic is identical.

The method first evaluates reporting framework, jurisdiction and dates; then
entity and transaction scope; contract and recognition facts; elections;
numerical capability; and account/context bindings. Every material finding
survives. A false required fact is ineligible, an unknown required fact is
unresolved, and an event outside the method is outside scope. A known false
finding must not disappear behind a later unknown one. Missing operative
authority never becomes source-supported permission.

The native operation \code{accounting.assess@1} remains read-only assessment.
Its \code{assess\_only} request and false posting permission cannot be promoted
by editing a flag. The separate \code{accounting.authorize@1} operation produces
a \code{PostingAuthorization} from the complete basis and a supported posting
profile. It permits posting only when the treatment is eligible, the method is
fully specified for the event and the branch mode matches the profile's
evidence scope. An unresolved or unsupported assessment is still a useful
completed assessment, but it contains no permission to post.

An authorized proposal contains both a journal and an independent economic
delta, each bound to the same basis and authorization. Journal lines name
account role, debit or credit, positive exact quantity, currency/context,
entity, book, event and method. Zero lines are omitted; a negative effect reverses
sides. The journal balances within each currency and book. The delta separately
checks quantities and contractual memory: inventory movements consume available
lots, cash follows settlement, share counts follow ownership events and debt
face amount remains distinct from carrying value.

The commit arguments have exactly four references:
\begin{lstlisting}
AccountingCommitArguments/v1 {
  basis_ref, authorization_ref, journal_ref, delta_ref
}
\end{lstlisting}
The expected opening state and branch generation appear once, in the outer
operation request. A reference to a native assessment in the authorization
field fails its reference type. False posting permission fails authorization.
An unknown posting schema returns unsupported operation with no journal commit.
All bindings are checked before looking up a prior receipt.

\begin{figure}[htbp]
\centering
\begin{tikzpicture}[node distance=5mm,>=Stealth,
 box/.style={draw,rounded corners,align=center,text width=29mm,minimum height=13mm,font=\small}]
\node[box] (basis) {Economic event\\and opening state};
\node[box,right=of basis] (auth) {Eligible method\\and permission};
\node[box,right=of auth] (proposal) {Journal and\\economic delta};
\node[box,right=of proposal] (commit) {Atomic state\\and receipt};
\draw[->] (basis)--(auth);
\draw[->] (auth)--(proposal);
\draw[->] (proposal)--(commit);
\end{tikzpicture}
\caption{Each step answers a different question. Accounting assessment may
inform the permission decision, but only the final transaction changes the book.}
\end{figure}

For a delivery, the basis identifies the exact performance obligation,
quantity, allocation and control finding. The opening state contains cumulative
revenue, billing, collection and inventory. The selected method derives the
recognition and lot movements; the report later consumes that same committed
path. Combining a journal from one method version with a delta or opening state
from another is invalid even if the totals happen to balance.

\section{A lost acknowledgement must not create a second sale}
\label{sec:replay}

Suppose a sale commits but the response is lost. The client knows the request,
but cannot tell whether the server committed it. Retrying under a new event ID
could duplicate revenue. Retrying the original request must instead return the
original immutable receipt.

The lookup key is
\[
  (\text{entity},\ \text{book},\ \text{branch},\ \text{business event}).
\]
Payload, event version, opening state and request settings belong to the complete
fingerprint, not this key. A renamed run or bumped event version cannot evade a
conflict. A separately commissioned hypothetical projection uses a new branch
namespace; reopening the same projection uses its existing one.

A correction of an already committed event is a new business event linked to
its predecessor and an admitted correction method, such as reversal and
replacement or a dated adjustment. It preserves the original receipt.
The accounting treatment must apply to the actual correction; absent treatment
authority leaves it unresolved rather than licensing a generic reversal.

Commit first validates input shape, closure, bindings and authorization in the
requested historical view. It then looks up the stable event key. An equal
fingerprint returns the old receipt with no additional commit; a different
fingerprint returns an event-identity conflict. Only an absent key proceeds to
the expected-state and generation comparison. This order is essential: after
a successful sale the current head has changed, so checking freshness before
replay would reject the very retry needed to resolve the lost response.

A new candidate is assessed and constructed without mutation. One transaction
then rechecks key absence and the expected head, inserts journal, delta, closing
state, dependencies and receipt, and advances the head. It commits all or none.
If another client has committed the identical event, return that receipt. A
different event that advanced the head makes this request stale. A changed
payload under the same key conflicts. Diagnostic or scratch records never count
as economic transitions.

A timeout consequently has three possible observations: committed, not committed
or unknown. Unknown retains the original key, fingerprint and expected head for
receipt lookup or exact retry. Refreshing the head changes the request and cannot
be presented as the same retry. Missing committed dependencies are an integrity
failure, not permission to reconstruct an approximately similar receipt.
